\begin{proof}[Proof of Lemma~\ref{lem:residual-variance-brownian}]
    Now~$R_1=\Var(\theta(t_1))=\sigma_0^2+\sigma^2t_1$ by definition.
    
    Suppose~$n>1$.
    The Brownian motion is a Gauss-Markov process, so we can use Lemma~\ref{lem:residual-variance-markov} to characterize~$R_n$.
    Define~$\beta_n$ and~$\gamma_n$ as in that lemma.
    Then
    %
    \begin{align*}
        \beta_n
        &= \frac{\sigma_0^2+\sigma^2\min\{t_n,t_{n-1}\}}{\sigma_0^2+\sigma^2t_{n-1}} \\
        &= 1
    \end{align*}
    %
    and
    %
    \begin{align*}
        \gamma_n
        &\overset{\star}{=} \Var(\theta(t_n))-\frac{\left(\Cov(\theta(t_n),\theta(t_{n-1}))\right)^2}{\Var(\theta(t_{n-1}))} \\
        &= \sigma_0^2+\sigma^2t_n-\frac{\left(\sigma_0^2+\sigma^2\min\{t_n,t_{n-1}\}\right)^2}{\sigma_0^2+\sigma^2t_{n-1}} \\
        &= \sigma^2(t_n-t_{n-1}),
    \end{align*}
    %
    where~$\star$ holds by Lemma~\ref{lem:multivariate-normal-conditional}.
    The result follows from Lemma~\ref{lem:residual-variance-markov}.
\end{proof}
